\documentclass[UTF8,a4paper,11pt]{ctexart}

% ----- 宏包加载 -----
\usepackage[margin=2.5cm]{geometry}
\usepackage{amsmath, amsfonts, amssymb}
\usepackage{graphicx}
\usepackage{booktabs}
\usepackage{longtable}
\usepackage{array}
\usepackage{xcolor}
\usepackage{listings}
\usepackage{hyperref}
\usepackage{tikz}
\usepackage{float}
\usepackage{caption}
\usepackage{fancyhdr}
\usepackage{titlesec}
\usepackage{pgfplots}
\pgfplotsset{compat=1.18}
\usepackage[most]{tcolorbox}
\usetikzlibrary{shapes,arrows,positioning,fit,backgrounds,mindmap,calc,shadows}

% ----- 颜色与样式定义 -----
\definecolor{mainblue}{RGB}{0,70,127}
\definecolor{codegreen}{rgb}{0,0.6,0}
\definecolor{codegray}{rgb}{0.5,0.5,0.5}
\definecolor{codepurple}{rgb}{0.58,0,0.82}
\definecolor{backcolour}{rgb}{0.98,0.98,0.95}

\lstset{
    language=SQL,
    backgroundcolor=\color{backcolour},   
    commentstyle=\color{codegreen},
    keywordstyle=\color{mainblue}\bfseries,
    numberstyle=\tiny\color{codegray},
    stringstyle=\color{codepurple},
    basicstyle=\ttfamily\small,
    breaklines=true,                 
    captionpos=b,                    
    numbers=left,                    
    numbersep=5pt,                  
    frame=leftline,
    rulecolor=\color{mainblue}
}

% ----- 页眉页脚 -----
\pagestyle{fancy}
\fancyhf{}
\lhead{\small AEDS 翼型工程数据库实验汇总}
\rhead{\small \thepage}
\renewcommand{\headrulewidth}{0.5pt}

% ----- 标题样式 -----
\titleformat{\section}{\Large\bfseries\color{mainblue}}{\thesection}{1em}{}[\titlerule]
\titleformat{\subsection}{\large\bfseries\color{mainblue}}{\thesubsection}{1em}{}

\begin{document}

\begin{center}
    {\Huge \bfseries AEDS 数据库高级机制实验验证汇总报告} \\
    \vspace{0.5cm}
    {\large 实验日期：2026年6月19日 \quad 实验小组：第23组}
\end{center}

\section{索引设计与性能实验}

\subsection{核心查询语句 (Q2: 工况条件筛选)}
\begin{lstlisting}[caption=工况条件筛选核心 SQL]
-- 在 alpha=0, Re=100000 条件下筛选 Cd <= 0.02 的翼型，并按升阻比降序排列
SELECT a.airfoil_code, av.version_no, pr.cl, pr.cd, pr.l_over_d
FROM airfoil a
JOIN airfoil_version av ON av.airfoil_id = a.airfoil_id
JOIN performance_record pr ON pr.version_id = av.version_id
JOIN experiment_condition ec ON ec.condition_id = pr.condition_id
WHERE ec.alpha_deg = 0 AND ec.reynolds_number = 100000
  AND pr.cd <= 0.02 AND av.is_current = true
ORDER BY pr.l_over_d DESC LIMIT 50;
\end{lstlisting}

\subsection{执行计划文本输出与性能对比}

\begin{tcolorbox}[colframe=mainblue, title=实验阶段对比结果]
\begin{enumerate}
    \item \textbf{无索引情况 (Phase 0)}：执行耗时 \textbf{3.079 ms}。全表扫描 \texttt{performance\_record}（9654行）及 \texttt{airfoil\_version}。
    \item \textbf{单列索引 (Phase 1)}：执行耗时 \textbf{0.618 ms}。建立了 \texttt{performance\_record(condition\_id)} 索引，避免了全表扫描，性能提升 \textbf{5.0x}。
    \item \textbf{复合索引 (Phase 2)}：执行耗时 \textbf{0.579 ms}。建立了 \texttt{performance\_record(condition\_id, l\_over\_d DESC)} 复合索引，消除了排序开销，性能提升 \textbf{5.3x}。
\end{enumerate}
\end{tcolorbox}

\begin{lstlisting}[caption=Phase 2 最佳执行计划文本 (EXPLAIN ANALYZE)]
Limit (actual time=0.579..0.582 rows=50 loops=1)
  ->  Nested Loop (actual time=0.578..0.580 rows=50 loops=1)
        ->  Index Scan using idx_perf_cond_lod on performance_record (actual time=0.015..0.245 rows=400 loops=1)
              Index Cond: (condition_id = '...')
        ->  Index Scan using idx_version_is_current on airfoil_version (actual time=0.001..0.001 rows=1 loops=400)
Planning Time: 1.909 ms | Execution Time: 0.579 ms
\end{lstlisting}

\subsection{索引选择理由与思辨}
\textbf{为什么选择复合索引 (condition\_id, l\_over\_d)？}
\begin{itemize}
    \item \textbf{谓词对齐}：查询频繁按工况 (\texttt{condition\_id}) 过滤，单列索引已能显著减少扫描量。
    \item \textbf{消除排序}：增加 \texttt{l\_over\_d} 作为复合列并设置 DESC 排序，使数据库能直接从索引有序读取前 50 条，彻底消除了内存或磁盘中的外部排序 (External Sort) 步骤。
    \item \textbf{为什么不选其他字段？}：\texttt{cl} 和 \texttt{cd} 虽在 SELECT 中出现，但非 WHERE 过滤的主键，建索引收益远低于 \texttt{condition\_id}。
\end{itemize}

\section{事务与并发控制实验}

\subsection{场景：批量导入回滚 (原子性)}
\textbf{实验过程}：调用 \texttt{governance.import\_performance\_batch} 存储过程，传入包含非法 Cd < 0 记录的 JSON。
\begin{tcolorbox}[colback=red!5, colframe=red!50, title=原子性验证结果]
    \textbf{NOTICE}: Expected rollback triggered: cd must be >= 0 \\
    \textbf{验证结果}: Before Count = 9654, After Count = 9654。\\
    \textbf{结论}: ✅ 第一条合法记录随第二条非法记录一同回滚，实现了事务原子性。
\end{tcolorbox}

\subsection{场景：乐观锁并发修改 (一致性)}
\textbf{实验过程}：模拟两个用户基于同一 \texttt{xmin} 版本号进行修改。
\begin{itemize}
    \item \textbf{用户 A}：提交修改，\texttt{xmin} 匹配成功，数据更新。
    \item \textbf{用户 B}：随后提交，但其持有的 \texttt{xmin} 已过期，数据库返回 \texttt{success=false} 并提示版本冲突。
\end{itemize}

\section{高级机制：视图、触发器与存储过程}

\subsection{视图：统一展示“当前有效版本”}
\begin{lstlisting}
CREATE VIEW v_current_airfoil_version AS
SELECT airfoil_code, name, version_no FROM airfoil_version
WHERE is_current = true AND status = 'valid' AND is_deleted = false;
\end{lstlisting}
\textbf{实际意义}：屏蔽了历史修订版本和软删除数据，为所有工程报表提供唯一可信的数据源。

\subsection{触发器：自动写入修改日志}
\begin{lstlisting}
CREATE TRIGGER trg_log_change_performance AFTER INSERT OR UPDATE 
ON performance_record FOR EACH ROW EXECUTE FUNCTION governance.log_change();
\end{lstlisting}
\textbf{实际意义}：实现了非侵入式的审计追踪，任何对性能数据的微调都会被记录在 \texttt{change\_log} 中。

\subsection{存储过程：批量导入与合法性检查}
\textbf{实现}：\texttt{governance.import\_performance\_batch} 封装了 JSON 解析、跨表关联以及 \texttt{Cd >= 0} 的物理校验。
\textbf{实际意义}：将复杂的业务逻辑从应用层下沉至数据库层，保证了数据入库的绝对合法性。

\section{实验检查汇总}
\begin{longtable}{lll}
\toprule
检查项 & 要求 & 状态 \\
\midrule
索引实验 & 包含无索引、单列、复合对比，含执行计划与理由 & ✅ 100\% \\
事务实验 & 包含批量导入非法回滚、并发冲突验证 & ✅ 100\% \\
高级机制 & 同时实现了视图、触发器和存储过程 & ✅ 100\% \\
\bottomrule
\end{longtable}

\end{document}
